% !TeX program = pdflatex
\documentclass[11pt,a4paper]{amsart}

\newcommand{\notelanguage}{finnish}
% Shared preamble for course notes and exercise sets.

\usepackage[T1]{fontenc}
\usepackage{lmodern}
\providecommand{\notelanguage}{english}
\usepackage[finnish,english]{babel}
\AtBeginDocument{\expandafter\selectlanguage\expandafter{\notelanguage}}

\newcommand{\notesfinnishlanguage}{finnish}
\ifx\notelanguage\notesfinnishlanguage
\newcommand{\notestheoremname}{Lause}
\newcommand{\noteslemmaname}{Lemma}
\newcommand{\notespropositionname}{Propositio}
\newcommand{\notescorollaryname}{Seurauslause}
\newcommand{\notesdefinitionname}{Määritelmä}
\newcommand{\notesexamplename}{Esimerkki}
\newcommand{\notesproblemname}{Tehtävä}
\newcommand{\notesexercisename}{Tehtävä}
\newcommand{\notesremarkname}{Huomautus}
\newcommand{\notessolutionname}{Ratkaisu}
\else
\newcommand{\notestheoremname}{Theorem}
\newcommand{\noteslemmaname}{Lemma}
\newcommand{\notespropositionname}{Proposition}
\newcommand{\notescorollaryname}{Corollary}
\newcommand{\notesdefinitionname}{Definition}
\newcommand{\notesexamplename}{Example}
\newcommand{\notesproblemname}{Problem}
\newcommand{\notesexercisename}{Exercise}
\newcommand{\notesremarkname}{Remark}
\newcommand{\notessolutionname}{Solution}
\fi

\usepackage[a4paper,margin=30mm]{geometry}
\usepackage{microtype}
\usepackage{graphicx}
\usepackage{enumitem}

\newcommand{\exerciseheading}[2]{%
  \par\noindent
  {\large\bfseries #1\par}
  \vspace{0.1em}
  \noindent{\large #2\par}
  \vspace{1.5em}
}

\usepackage{amsmath}
\usepackage{amssymb}
\usepackage{amsthm}
\usepackage{mathtools}
\usepackage{aliascnt}

% Linear systems: the argument gives the number of alignment pairs.
\newenvironment{linearsystem}[1]
  {\left\{\begin{alignedat}{#1}}
  {\end{alignedat}\right.}

\usepackage[hidelinks,pdfusetitle]{hyperref}

\numberwithin{equation}{section}
\setlist{itemsep=0.25em,topsep=0.5em}

\theoremstyle{plain}
\newtheorem{theorem}{\notestheoremname}[section]
\newaliascnt{lemma}{theorem}
\newtheorem{lemma}[lemma]{\noteslemmaname}
\aliascntresetthe{lemma}
\newaliascnt{proposition}{theorem}
\newtheorem{proposition}[proposition]{\notespropositionname}
\aliascntresetthe{proposition}
\newaliascnt{corollary}{theorem}
\newtheorem{corollary}[corollary]{\notescorollaryname}
\aliascntresetthe{corollary}

\theoremstyle{definition}
\newaliascnt{definition}{theorem}
\newtheorem{definition}[definition]{\notesdefinitionname}
\aliascntresetthe{definition}
\newaliascnt{example}{theorem}
\newtheorem{example}[example]{\notesexamplename}
\aliascntresetthe{example}
\newaliascnt{problem}{theorem}
\newtheorem{problem}[problem]{\notesproblemname}
\aliascntresetthe{problem}
\newtheorem{exercise}{\notesexercisename}

\theoremstyle{remark}
\newaliascnt{remark}{theorem}
\newtheorem{remark}[remark]{\notesremarkname}
\aliascntresetthe{remark}

\newenvironment{solution}
{
\begin{proof}[\notessolutionname]}
  {
\end{proof}}

\usepackage[nameinlink,noabbrev]{cleveref}

\crefname{theorem}{theorem}{theorems}
\crefname{lemma}{lemma}{lemmas}
\crefname{proposition}{proposition}{propositions}
\crefname{corollary}{corollary}{corollaries}
\crefname{definition}{definition}{definitions}
\crefname{example}{example}{examples}
\crefname{problem}{problem}{problems}
\crefname{exercise}{exercise}{exercises}
\crefname{remark}{remark}{remarks}

\ifx\notelanguage\notesfinnishlanguage
\crefname{theorem}{lause}{lauseet}
\crefname{lemma}{lemma}{lemmat}
\crefname{proposition}{propositio}{propositiot}
\crefname{corollary}{seurauslause}{seurauslauseet}
\crefname{definition}{määritelmä}{määritelmät}
\crefname{example}{esimerkki}{esimerkit}
\crefname{problem}{tehtävä}{tehtävät}
\crefname{exercise}{tehtävä}{tehtävät}
\crefname{remark}{huomautus}{huomautukset}
\fi

\newcommand{\NN}{\mathbb{N}}
\newcommand{\NNone}{\mathbb{N}_1}
\newcommand{\ZZ}{\mathbb{Z}}
\newcommand{\QQ}{\mathbb{Q}}
\newcommand{\RR}{\mathbb{R}}
\newcommand{\CC}{\mathbb{C}}
\newcommand{\PP}{\mathbb{P}}

\DeclareMathOperator{\im}{im}
\DeclareMathOperator{\rank}{rank}
\DeclareMathOperator{\syt}{syt}
\DeclarePairedDelimiter{\abs}{\lvert}{\rvert}
\DeclarePairedDelimiter{\norm}{\lVert}{\rVert}
\newcommand{\set}[1]{\left\lbrace#1\right\rbrace}
\newcommand{\maxset}[1]{\max\set{#1}}
\newcommand{\minset}[1]{\min\set{#1}}
\newcommand{\powerset}[1]{\mathcal{P}\!\left(#1\right)}


\title{Lineaarialgebra ja matriisilaskenta I}
\author{Joona Piirainen}
\date{Syksy 2026 -- periodi 2}

\begin{document}

\begin{abstract}
  Luentomuistiinpanot Helsingin yliopiston kurssilta
  Lineaarialgebra ja matriisilaskenta I (MAT11002).
  Kurssi järjestetään aikavälillä 27.10.2026--14.12.2026.
\end{abstract}

\maketitle

\tableofcontents

\section{Kurssitiedot}

Kurssin viralliset tiedot löytyvät
\href{https://studies.helsinki.fi/kurssit/opintojakso/otm-0a7d069a-f394-4a4a-8009-ab28a38e4419/MAT11002}{Helsingin
yliopiston kurssisivulta}.

\section{Muistiinpanot}

\subsection{27.10.2026: lineaariset yhtälöryhmät}

Lineaarisen yhtälöryhmän ratkaisu toteuttaa kaikki ryhmän yhtälöt
samanaikaisesti. Geometrisesti ratkaiseminen tarkoittaa yhtälöiden
ratkaisujoukkojen yhteisten pisteiden etsimistä.

\subsubsection*{Suorat tasossa ja niiden leikkauspisteet}

Yhtälö $y=ax+b$, missä $a,b\in\RR$, esittää tason suoraa.
Suora voidaan ajatella pistejoukkona
\[
  L=\{(x,y)\in\RR^2:y=ax+b\}.
\]
Kahden suoran $y=ax+b$ ja $y=a'x+b'$ yhteiset pisteet muodostavat
joukon $L\cap L'$. Ne saadaan ratkaisemalla molemmat yhtälöt yhtä aikaa.

\begin{example}[Kahden suoran leikkaus]
  Ratkaistaan yhtälöryhmä
  \[
    \begin{cases}
      y=2x+3,\\
      y=-x+2.
    \end{cases}
  \]
  Koska molemmat oikeat puolet ovat yhtä suuret kuin $y$, saadaan
  \[
    2x+3=-x+2
    \quad\Longleftrightarrow\quad
    3x=-1
    \quad\Longleftrightarrow\quad
    x=-\frac13.
  \]
  Sijoittamalla toiseen yhtälöön saadaan
  \[
    y=-\left(-\frac13\right)+2=\frac73.
  \]
  Suorien ainoa leikkauspiste on siis
  \[
    \boxed{(x,y)=\left(-\frac13,\frac73\right).}
  \]
\end{example}

\newpage
\subsubsection*{Kahden suoran yhteiset pisteet}
\begin{center}
  \renewcommand{\arraystretch}{1.35}
  \begin{tabular}{lll}
    \hline
    Ehto & Suorien sijainti & Ratkaisuja\\
    \hline
    $a\ne a'$ & Leikkaavat yhdessä pisteessä & Yksi\\
    $a=a'$, $b\ne b'$ & Erilliset yhdensuuntaiset suorat & Ei yhtään\\
    $a=a'$, $b=b'$ & Sama suora & Äärettömän monta\\
    \hline
  \end{tabular}
\end{center}

\subsubsection*{Suoran ja tason yleinen yhtälö}

Muoto $y=ax+b$ ei kata pystysuoria suoria $x=c$. Kaikki tason suorat
voidaan esittää muodossa
\[
  \boxed{a_1x_1+a_2x_2=b,\qquad (a_1,a_2)\ne(0,0),}
\]
missä $a_1,a_2,b\in\RR$. Esimerkiksi $x_1=c$ saadaan valitsemalla
$a_1=1$, $a_2=0$ ja $b=c$. Suoran pistejoukko on
\[
  L=\{(x_1,x_2)\in\RR^2:a_1x_1+a_2x_2=b\}.
\]

\begin{remark}[Yhtälön kertominen vakiolla]
  Kun $\lambda\ne0$, yhtälön molemmat puolet voidaan kertoa luvulla
  $\lambda$ muuttamatta ratkaisujoukkoa:
  \[
    a_1x_1+a_2x_2=b
    \quad\Longleftrightarrow\quad
    \lambda a_1x_1+\lambda a_2x_2=\lambda b.
  \]
\end{remark}

Vastaavasti avaruuden $\RR^3$ tason yleinen yhtälö on
\[
  a_1x_1+a_2x_2+a_3x_3=b,
  \qquad (a_1,a_2,a_3)\ne(0,0,0).
\]

\subsubsection*{Yhtälöryhmä, ratkaisu ja ratkaisujoukko}

\begin{definition}[Lineaarinen yhtälöryhmä]
  Olkoot $m,n$ positiivisia kokonaislukuja ja
  $a_{ij},b_i\in\RR$, missä $i=1,\ldots,m$ ja $j=1,\ldots,n$.
  Ryhmää
  \begin{equation}\label{eq:linear-system}
    \left\{
      \begin{aligned}
        a_{11}x_1+\cdots+a_{1n}x_n&=b_1,\\
        &\ \vdots\\
        a_{m1}x_1+\cdots+a_{mn}x_n&=b_m
      \end{aligned}\right.
      \tag{YR}
    \end{equation}
    kutsutaan \emph{lineaariseksi yhtälöryhmäksi}.
    Siinä on $m$ yhtälöä ja $n$ tuntematonta.
    Kerroin $a_{ij}$ on yhtälön $i$ muuttujan $x_j$ kerroin.
  \end{definition}

  \begin{definition}[Ratkaisu ja ratkaisujoukko]
    Vektori $(x_1,\ldots,x_n)\in\RR^n$ on ryhmän
    \eqref{eq:linear-system} \emph{ratkaisu}, jos se toteuttaa ryhmän
    kaikki yhtälöt. Kaikkien ratkaisujen joukko on
    \[
      R=\left\{(x_1,\ldots,x_n)\in\RR^n:
      \sum_{j=1}^n a_{ij}x_j=b_i\ \text{kaikilla }i=1,\ldots,m\right\}.
    \]
  \end{definition}

  \noindent\textbf{Geometrinen tulkinta.}
  Jos $L_i$ on yhtälön $i$ ratkaisujoukko, koko ryhmän ratkaisujoukko on
  $R=L_1\cap\cdots\cap L_m$.

  \newpage
  \subsubsection*{Avaruuden \texorpdfstring{$\RR^n$}{Rⁿ} laskutoimitukset}

  Avaruuden $\RR^n$ alkiot ovat vektoreita $(x_1,\ldots,x_n)$.
  Vektorien yhteenlasku ja skalaarilla kertominen tehdään
  \emph{koordinaateittain}.

  \begin{definition}[Yhteenlasku ja skalaarilla kertominen]
    Yhteenlasku $+\colon\RR^n\times\RR^n\to\RR^n$ määritellään kaavalla
    \[
      \boxed{(x_1,\ldots,x_n)+(y_1,\ldots,y_n)
      =(x_1+y_1,\ldots,x_n+y_n).}
    \]
    Skalaarilla kertominen $\RR\times\RR^n\to\RR^n$ määritellään kaavalla
    \[
      \boxed{c(x_1,\ldots,x_n)=(cx_1,\ldots,cx_n).}
    \]
    Tässä $c\in\RR$ on \emph{skalaari}.
  \end{definition}

  \begin{example}
    Kun $u=(1,-2)$ ja $v=(3,4)$, saadaan
    \[
      u+v=(4,2),\qquad 2u=(2,-4).
    \]
  \end{example}

  \subsubsection*{Vektorien laskusäännöt}
  Kaikilla $u,v,w\in\RR^n$ ja $c,d\in\RR$ pätevät
  \begin{align*}
    u+v&=v+u, & (u+v)+w&=u+(v+w),\\
    u+0&=u, & u+(-u)&=0,\\
    c(u+v)&=cu+cv, & (c+d)u&=cu+du,\\
    c(du)&=(cd)u, & 1u&=u.
  \end{align*}
  Näiden laskutoimitusten kanssa $\RR^n$ on vektoriavaruus kunnan
  $\RR$ yli. Yleisemmin puhutaan vektoriavaruudesta kunnan $K$ yli;
  renkaan yli vastaava käsite on $K$-moduuli.

  \subsubsection*{Yhtälöt $0=0$ ja $0=b_i$}

  Yhtälöryhmässä voi esiintyä rivi, jonka kaikki kertoimet ovat nollia:
  \[
    0x_1+0x_2+\cdots+0x_n=b_i.
  \]
  Tällöin on erotettava kaksi tapausta.
  \begin{description}[style=nextline,leftmargin=0pt]
    \item[$b_i=0$: yhtälö $0=0$.]
      Rivi toteutuu kaikilla vektoreilla eikä aseta ratkaisuille mitään
      ehtoa. Sen voi poistaa muuttamatta ryhmän ratkaisujoukkoa.
    \item[$b_i\ne0$: yhtälö $0=b_i$.]
      Rivi ei toteudu millään vektorilla. Siksi koko yhtälöryhmällä
      ei ole ratkaisuja: $\boxed{R=\varnothing}$.
  \end{description}

  \end{document}
